% Problem 5 — Slice Filtration for N_infty Operads (Blumberg)
% Solution v6 — claude-opus-4-7 — April 30, 2026
\documentclass[12pt]{article}
\usepackage{fullpage}
\usepackage[T1]{fontenc}
\usepackage[utf8]{inputenc}
\usepackage{lmodern}
\usepackage{microtype}
\usepackage{amsmath,amssymb,amsthm,amsfonts}
\usepackage{hyperref}

\newtheorem{theorem}{Theorem}
\newtheorem{lemma}[theorem]{Lemma}
\newtheorem{proposition}[theorem]{Proposition}
\newtheorem{corollary}[theorem]{Corollary}
\theoremstyle{definition}
\newtheorem{definition}[theorem]{Definition}
\newtheorem{example}[theorem]{Example}
\newtheorem{remark}[theorem]{Remark}

\newcommand{\Sp}{\mathrm{Sp}}
\newcommand{\SpG}{\mathrm{Sp}^{G}}
\newcommand{\SpH}{\mathrm{Sp}^{H}}
\newcommand{\SpK}{\mathrm{Sp}^{K}}
\newcommand{\Spec}{\mathrm{Sp}}
\newcommand{\Phii}[1]{\Phi^{#1}}
\newcommand{\PhiH}{\Phi^{H}}
\newcommand{\PhiK}{\Phi^{K}}
\newcommand{\PhiG}{\Phi^{G}}
\newcommand{\OO}{\mathcal{O}}
\newcommand{\FF}{\mathcal{F}}
\newcommand{\NN}{\mathbb{N}}
\newcommand{\ZZ}{\mathbb{Z}}
\newcommand{\Sm}{\mathbb{S}}
\newcommand{\admis}{\mathrm{Adm}}
\newcommand{\fix}[1]{\widetilde{E\mathcal{F}[#1]}}
\newcommand{\Esub}[1]{E\mathcal{F}[#1]}

\title{Problem 5: Slice Filtration for $N_\infty$ Operads\\
A Geometric Fixed-Point Characterisation}
\author{}
\date{}

\begin{document}
\maketitle

\begin{abstract}
Fix a finite group $G$ and an incomplete transfer system $\OO$ on $G$ in the
sense of Blumberg--Hill. We define the $\OO$-adapted slice filtration on the
$G$-equivariant stable homotopy category $\SpG$ following the
Hill--Hopkins--Ravenel template, but using only those equivariant cells
$G_{+}\wedge_{H}S^{k\rho_{H}}$ for which $H$ is $\OO$-admissible. The main
theorem (Theorem~\ref{thm:main}) states and proves: \emph{a connective $G$-spectrum
$X$ lies in the $\OO$-slice connective subcategory $\tau_{\geq n}^{\OO}\SpG$
if and only if for every $\OO$-admissible subgroup $H\leq G$ the geometric
fixed-point spectrum $\PhiH(X)$ is $(\lfloor n/|H|\rfloor-1)$-connected.}
The proof has two halves. The forward direction verifies the connectivity
claim on each generating cell by an explicit double-coset computation of the
geometric fixed points. The reverse direction constructs $X$ as an
$\OO$-cellular homotopy colimit, using the standard isotropy-separation
cofibre sequence to detect cell-mapping spaces by the geometric fixed points
in the connective range.
\end{abstract}

\tableofcontents
\bigskip

%-------------------------------------------------------------
\section{Conventions and notation}
%-------------------------------------------------------------

We work in the genuine $G$-equivariant stable homotopy category $\SpG$ for a
finite group $G$ (e.g., the homotopy category of orthogonal $G$-spectra
inverted at all representation spheres). For $H\leq G$ we use:

\begin{itemize}
\item $i_{H}^{\ast}\colon \SpG\to\SpH$: restriction; left adjoint
$G_{+}\wedge_{H}(-)$, right adjoint $F_{H}(G_{+},-)$. Wirthmüller equivalence:
$G_{+}\wedge_{H}Y\simeq F_{H}(G_{+},Y)$.
\item $\rho_{H}$: the regular real $H$-representation, of dimension $|H|$.
\item $S^{V}$: the one-point compactification of an orthogonal
$H$-representation $V$, viewed as $H$-spectrum.
\item $\PhiH\colon \SpG\to\Spec$: the geometric fixed-point functor at $H$.
By definition $\PhiH(X)\simeq(\fix{H}\wedge i_{H}^{\ast}X)^{H}$, where
$\mathcal{F}[H]$ is the family of \emph{proper} subgroups of $H$, $\Esub{H}$ is
the universal $H$-CW-complex with isotropy in $\mathcal{F}[H]$, and $\fix{H}$
is the cofibre of $\Esub{H}_{+}\to S^{0}$. Writing the source as $\SpG$ uses
$i_{H}^{\ast}$ implicitly.
\item $\pi_{n}^{H}(X)=[S^{n},X]_{H}=[S^{n},i_{H}^{\ast}X]_{H}$: the genuine
$H$-equivariant homotopy.
\item A $G$-spectrum $X$ is \emph{connective} if $\pi_{k}^{H}(X)=0$ for every
$H\leq G$ and every $k<0$.
\item For a nonequivariant spectrum $E$, ``$E$ is $m$-connected'' means
$\pi_{k}E=0$ for $k\leq m$. ``$(-1)$-connected'' means connective.
\end{itemize}

We will use the following two standard facts.

\begin{lemma}[Wirthmüller adjunction]\label{lem:wirth}
For any $H\leq G$, any $H$-spectrum $Y$, any $G$-spectrum $X$:
\begin{equation}
\bigl[G_{+}\wedge_{H}Y,\,X\bigr]_{G}\;\cong\;\bigl[Y,\,i_{H}^{\ast}X\bigr]_{H},
\qquad
\bigl[i_{H}^{\ast}X,\,Y\bigr]_{H}\;\cong\;\bigl[X,\,G_{+}\wedge_{H}Y\bigr]_{G}.
\label{eq:wirth}
\end{equation}
\end{lemma}

\begin{proof}
Standard for genuine equivariant spectra; Mandell--May, \emph{Equivariant
orthogonal spectra and $S$-modules}, Theorem V.2.2.
\end{proof}

\begin{lemma}[Geometric fixed points: induction and spheres]\label{lem:dcp}
For $H,K\leq G$ and any $H$-spectrum $Y$,
\begin{equation}
\PhiK\bigl(G_{+}\wedge_{H}Y\bigr)
\;\simeq\;\bigvee_{[g]\in K\backslash G/H}\Phi^{K\cap gHg^{-1}}\bigl(c_{g}Y\bigr),
\label{eq:dcp}
\end{equation}
where $c_{g}Y$ is $Y$ with the $gHg^{-1}$-action $g'\cdot y:=g^{-1}g'g\cdot y$.
For any orthogonal $H$-representation $V$ and any $J\leq H$,
\begin{equation}
\Phi^{J}(S^{V})\;\simeq\;S^{V^{J}}.
\label{eq:phi-sphere}
\end{equation}
\end{lemma}

\begin{proof}
Hill--Hopkins--Ravenel, \emph{On the non-existence of elements of Kervaire
invariant one}, Annals 184 (2016), Proposition 2.59 and Proposition 2.55.
\end{proof}

We will use repeatedly the dimension count for fixed points of the regular
representation:

\begin{lemma}[Fixed-points of $\rho_{H}$]\label{lem:rho-fix}
For any $K\leq H$, $\dim(\rho_{H})^{K}=|H|/|K|=[H:K]$.
\end{lemma}

\begin{proof}
The restriction $\rho_{H}\downarrow_{K}$ of the regular $H$-representation to
$K$ is isomorphic to $(\rho_{K})^{[H:K]}$ (the regular $K$-representation
$[H:K]$ times). Indeed, $\rho_{H}=\mathbb{R}[H]$ as a left $H$-module;
restricting to $K$ partitions $H$ into $[H:K]$ left $K$-cosets, each yielding a
copy of $\mathbb{R}[K]=\rho_{K}$. Thus
\[
(\rho_{H})^{K}=\bigl((\rho_{K})^{[H:K]}\bigr)^{K}=
((\rho_{K})^{K})^{[H:K]}.
\]
Now $(\rho_{K})^{K}$ is one-dimensional, spanned by $\sum_{k\in K}k\in
\mathbb{R}[K]$. Hence $\dim(\rho_{H})^{K}=[H:K]=|H|/|K|$.
\end{proof}

%-------------------------------------------------------------
\section{Transfer systems and $N_\infty$ operads}\label{sec:trans}
%-------------------------------------------------------------

\begin{definition}[Transfer system]\label{def:trans}
A \emph{transfer system} on $G$ is a collection
$\OO\subseteq\{(H,K):H\leq K\leq G\}$ of pairs of subgroups, written
$H\to_{\OO}K$, satisfying:
\begin{itemize}
\item[\textup{(T1)}] \emph{Reflexivity:} $H\to_{\OO}H$ for every $H\leq G$.
\item[\textup{(T2)}] \emph{Conjugation closure:} if $H\to_{\OO}K$, then for
every $g\in G$, $gHg^{-1}\to_{\OO}gKg^{-1}$.
\item[\textup{(T3)}] \emph{Restriction closure:} if $H\to_{\OO}K$ and
$L\leq K$, then $H\cap L\to_{\OO}L$.
\item[\textup{(T4)}] \emph{Transitivity:} if $H\to_{\OO}K$ and $K\to_{\OO}L$,
then $H\to_{\OO}L$.
\end{itemize}
The transfer system is \emph{complete} if $H\to_{\OO}K$ for every
$H\leq K\leq G$, and \emph{incomplete} otherwise.
\end{definition}

\begin{definition}[$\OO$-admissible subgroup]\label{def:adm}
A subgroup $H\leq G$ is \emph{$\OO$-admissible} if $\{e\}\to_{\OO}H$. We write
\[
\admis(\OO)\;:=\;\{H\leq G\,:\,H\text{ is }\OO\text{-admissible}\}.
\]
\end{definition}

\begin{lemma}[$\admis(\OO)$ is a family]\label{lem:adm-family}
$\admis(\OO)$ is closed under conjugation and under taking subgroups: it is a
family of subgroups of $G$.
\end{lemma}

\begin{proof}
\emph{Conjugation:} if $H\in\admis(\OO)$ then $\{e\}\to_{\OO}H$, and (T2)
gives $g\{e\}g^{-1}\to_{\OO}gHg^{-1}$, i.e., $\{e\}\to_{\OO}gHg^{-1}$.

\emph{Subgroups:} if $K\leq H$ and $H\in\admis(\OO)$, then
$\{e\}\to_{\OO}H$. By (T3) with $L:=K\leq H$,
\[
\{e\}\cap K\;\to_{\OO}\;K,
\quad\text{i.e.,}\quad \{e\}\to_{\OO}K,
\]
so $K\in\admis(\OO)$. Note $\{e\}\in\admis(\OO)$ by (T1).
\end{proof}

\begin{definition}[$N_\infty$ operad]\label{def:Ninfty}
An \emph{$N_\infty$ operad} for $G$ is a $G$-operad $\mathcal{C}$ such that
each $\mathcal{C}(n)$ is $\Sigma_{n}$-free and the fixed points
$\mathcal{C}(n)^{\Gamma}$ are either empty or contractible for every
$\Gamma\leq G\times\Sigma_{n}$. The Blumberg--Hill correspondence (\emph{Adv.
Math.} 285 (2015), Theorem 1.4) gives a bijection
\[
\{N_\infty\text{ operads on }G\}/\simeq
\;\;\longleftrightarrow\;\;
\{\text{transfer systems on }G\}.
\]
A transfer $H\to_{\OO}K$ corresponds to the existence of a norm
$N_{H}^{K}$ in the operad. We use only the transfer system $\OO$ and the
admissible family $\admis(\OO)$.
\end{definition}

\begin{example}\label{ex:trans-systems}
\begin{itemize}
\item Trivial transfer system $\OO_{\mathrm{tr}}=\{(H,H):H\leq G\}$:
$\admis(\OO_{\mathrm{tr}})=\{\{e\}\}$, corresponding to a naive $E_{\infty}$
operad.
\item Complete transfer system $\OO_{\mathrm{co}}=\{(H,K):H\leq K\leq G\}$:
$\admis(\OO_{\mathrm{co}})=\{H:H\leq G\}$, corresponding to a complete
$G$-$E_{\infty}$ operad.
\item For $G=C_{p^{2}}$ with chain $\{e\}<C_{p}<C_{p^{2}}$, the transfer
system
\[\OO=\{(\{e\},\{e\}),\,(C_{p},C_{p}),\,(C_{p^{2}},C_{p^{2}}),\,(\{e\},C_{p})\}\]
satisfies (T1)--(T4) and yields $\admis(\OO)=\{\{e\},C_{p}\}$; this is genuinely
incomplete.
\end{itemize}
\end{example}

%-------------------------------------------------------------
\section{Slice cells and the $\OO$-adapted slice filtration}\label{sec:filt}
%-------------------------------------------------------------

Fix a transfer system $\OO$ on $G$ throughout this section.

\begin{definition}[Slice cell]\label{def:cell}
For $H\leq G$ and $k\in\ZZ$, the \emph{slice cell of type $(H,k)$} is
\[
\widehat{S}(H,k)\;:=\;G_{+}\wedge_{H}S^{k\rho_{H}}\;\in\;\SpG.
\]
Its \emph{slice dimension} is $\dim\widehat{S}(H,k):=k\dim\rho_{H}=k|H|$.
\end{definition}

\begin{definition}[$\OO$-cells and the $\OO$-slice filtration]\label{def:tauO}
The \emph{$\OO$-slice cells} are the slice cells with $\OO$-admissible source:
\[
\mathcal{C}_{\OO}\;:=\;\bigl\{\widehat{S}(H,k):H\in\admis(\OO),\,k\in\ZZ\bigr\}.
\]
For $n\in\ZZ$, define $\tau_{\geq n}^{\OO}\SpG$ to be the smallest full
subcategory of $\SpG$ containing every $\widehat{S}(H,k)$ with
$H\in\admis(\OO)$ and $k|H|\geq n$, and closed under arbitrary small homotopy
colimits and extensions. Equivalently, $\tau_{\geq n}^{\OO}\SpG$ is the
localising subcategory generated by the set
$\{\widehat{S}(H,k):H\in\admis(\OO),\,k|H|\geq n\}$. We say $X\in\SpG$ is
\emph{$\OO$-$n$-slice-connective} if $X\in\tau_{\geq n}^{\OO}\SpG$.
\end{definition}

\begin{remark}\label{rem:HHR-special}
The case $\OO=\OO_{\mathrm{co}}$ recovers the Hill--Hopkins--Ravenel
filtration: $\tau_{\geq n}^{\OO_{\mathrm{co}}}\SpG=\tau_{\geq
n}^{\mathrm{HHR}}\SpG$. For incomplete $\OO$, the filtration is coarser
($\OO\subseteq\OO_{\mathrm{co}}$ implies $\tau_{\geq
n}^{\OO}\SpG\subseteq\tau_{\geq n}^{\mathrm{HHR}}\SpG$, since fewer cells
generate a smaller localising subcategory).
\end{remark}

%-------------------------------------------------------------
\section{Geometric fixed points of slice cells}\label{sec:cells-GF}
%-------------------------------------------------------------

The forward direction of the main theorem rests on a clean computation of
$\PhiK(\widehat{S}(H,k))$.

\begin{lemma}[Geometric fixed points of slice cells]\label{lem:GF-cells}
For all $H,K\leq G$ and all $k\in\ZZ$,
\begin{equation}
\PhiK\bigl(\widehat{S}(H,k)\bigr)
\;\simeq\;\bigvee_{[g]\in K\backslash G/H}S^{\,k\,|gHg^{-1}|/|K\cap gHg^{-1}|}.
\label{eq:cell-GF}
\end{equation}
In particular, when $k\geq 0$, $\PhiK(\widehat{S}(H,k))$ is a wedge of
non-negative-dimensional spheres, and the connectivity of this wedge equals
\[
\Bigl(\min_{[g]}k\,|gHg^{-1}|/|K\cap gHg^{-1}|\Bigr)-1
\;\;\geq\;\;
k\,|H|/|K|\,-\,1.
\]
\end{lemma}

\begin{proof}
Apply Lemma~\ref{lem:dcp} to $Y=S^{k\rho_{H}}$:
\[
\PhiK\bigl(G_{+}\wedge_{H}S^{k\rho_{H}}\bigr)
\;\simeq\;\bigvee_{[g]\in K\backslash G/H}
\Phi^{K\cap gHg^{-1}}\bigl(c_{g}S^{k\rho_{H}}\bigr).
\]
For each $g$, $c_{g}S^{k\rho_{H}}=S^{k\rho_{gHg^{-1}}}$ (the regular
representation transports to the regular representation under conjugation,
since $c_{g}\rho_{H}\cong\rho_{gHg^{-1}}$ as orthogonal representations).
Setting $J:=K\cap gHg^{-1}\leq gHg^{-1}$, by \eqref{eq:phi-sphere} of
Lemma~\ref{lem:dcp},
\[
\Phi^{J}\bigl(S^{k\rho_{gHg^{-1}}}\bigr)\;\simeq\;
S^{k\dim(\rho_{gHg^{-1}})^{J}}.
\]
By Lemma~\ref{lem:rho-fix} applied to $J\leq gHg^{-1}$,
$\dim(\rho_{gHg^{-1}})^{J}=|gHg^{-1}|/|J|=|gHg^{-1}|/|K\cap gHg^{-1}|$. This
proves \eqref{eq:cell-GF}.

For $k\geq 0$ each exponent is non-negative, so each summand is a
non-negatively-dimensioned sphere; the wedge is $(m-1)$-connected where $m$ is
the minimum exponent. Since $|K\cap gHg^{-1}|\leq|K|$,
\[
k\,|gHg^{-1}|/|K\cap gHg^{-1}|=k\,|H|/|K\cap gHg^{-1}|
\;\geq\;k\,|H|/|K|.
\]
The minimum is therefore at least $k|H|/|K|$, so the wedge is at least
$(k|H|/|K|-1)$-connected.
\end{proof}

\begin{corollary}[Cell connectivity bound]\label{cor:cell-conn}
Let $H,K\leq G$, $k\geq 0$, and set $n:=k|H|$. Then
$\PhiK(\widehat{S}(H,k))$ is $(\lfloor n/|K|\rfloor-1)$-connected.
\end{corollary}

\begin{proof}
By Lemma~\ref{lem:GF-cells}, $\PhiK(\widehat{S}(H,k))$ is at least
$(k|H|/|K|-1)$-connected, i.e., $(n/|K|-1)$-connected. Each exponent
$k|H|/|K\cap gHg^{-1}|$ in the wedge is a non-negative integer (since
$|K\cap gHg^{-1}|$ divides $|H|$, the quotient $|H|/|K\cap gHg^{-1}|$ is a
positive integer when $H$ is finite), so the minimum exponent is an integer
$\geq k|H|/|K|=n/|K|$, hence $\geq\lceil n/|K|\rceil\geq\lfloor n/|K|\rfloor$.
The connectivity bound $(\lfloor n/|K|\rfloor-1)$ follows.
\end{proof}

\begin{remark}\label{rem:sharp}
The bound in Corollary~\ref{cor:cell-conn} is sharp at $K=H$: for $g=e$ the
wedge summand has exponent exactly $k|H|/|H|=k$, so $\PhiH(\widehat{S}(H,k))$
contains a copy of $S^{k}$, hence is exactly $k$-connective and
$(k-1)$-connected. The full computation:
\[
\PhiH\bigl(\widehat{S}(H,k)\bigr)\simeq\bigvee_{[g]\in H\backslash G/H}
S^{\,k\,|H|/|H\cap gHg^{-1}|}.
\]
For $G=H$, $H\backslash G/H=\{[e]\}$ and the formula reduces to
$\PhiG(\widehat{S}(G,k))\simeq S^{k}$, recovering the standard fact
$\PhiG(S^{k\rho_{G}})=S^{k}$.
\end{remark}

%-------------------------------------------------------------
\section{The cell-mapping lemma in the connective range}\label{sec:cell-map}
%-------------------------------------------------------------

The reverse direction of the main theorem requires that maps out of an
$\OO$-cell into a connective spectrum $X$ be detected by $\PhiH X$.

\begin{lemma}[Isotropy-separation cofibre sequence]\label{lem:isotrop}
For any $H$-spectrum $Z$, there is a natural cofibre sequence in $\SpH$:
\begin{equation}
\Esub{H}_{+}\wedge Z\;\to\;Z\;\to\;\fix{H}\wedge Z.
\label{eq:isotrop}
\end{equation}
Applying the categorical fixed-point functor $(-)^{H}\colon\SpH\to\Spec$
yields a cofibre sequence in $\Spec$:
\begin{equation}
\bigl(\Esub{H}_{+}\wedge Z\bigr)^{H}\;\to\;Z^{H}\;\to\;\PhiH(Z),
\label{eq:fix-cof}
\end{equation}
the right-hand identification being the definition of $\PhiH$.
\end{lemma}

\begin{proof}
Standard isotropy separation; HHR §B.10.
\end{proof}

\begin{lemma}[Borel term connectivity]\label{lem:Borel-conn}
Let $H\leq G$ with $|H|\geq 2$, let $Y\in\SpH$ be connective, and let
$j\geq 0$. Let $p_{H}$ denote the smallest prime dividing $|H|$. Then the
``Borel term''
\begin{equation}
B_{j}(Y)\;:=\;\bigl(\Esub{H}_{+}\wedge S^{-j\rho_{H}}\wedge Y\bigr)^{H}
\;\in\;\Spec
\label{eq:Borel-def}
\end{equation}
is $(jp_{H}-1)$-connected; in particular it is $(j-1)$-connected
\textup{(}since $p_{H}\geq 2$\textup{)}. When $H=\{e\}$, $\Esub{\{e\}}=\emptyset$ and
$B_{j}(Y)=\ast$.
\end{lemma}

\begin{proof}
The case $H=\{e\}$ is immediate since $\mathcal{F}[\{e\}]=\emptyset$ and the
universal space $\Esub{\{e\}}$ is contractible (or empty in pointed
spectra). For $|H|\geq 2$, filter $\Esub{H}$ by its $H$-CW skeleta. By
definition of $\Esub{H}$ as the universal $H$-CW complex with isotropy in
$\mathcal{F}[H]$ (proper subgroups of $H$), every cell has the form
$H/L_{+}\wedge D^{m}$ with $L\lneq H$ proper and $m\geq 0$.

The $m$-th filtration quotient of $\Esub{H}_{+}\wedge S^{-j\rho_{H}}\wedge Y$
is a wedge over cells of the form
\[
H/L_{+}\wedge S^{m}\wedge S^{-j\rho_{H}}\wedge Y
\;\simeq\;H_{+}\wedge_{L}\bigl(S^{m}\wedge S^{-j\rho_{H}}\wedge i_{L}^{\ast}Y\bigr).
\]
By the categorical-fixed-points form of the Wirthmüller isomorphism
(Lemma~\ref{lem:wirth}; cf.\ HHR §B.4), for $L\lneq H$ and $W\in\SpL$,
\begin{equation}
\bigl(H_{+}\wedge_{L}W\bigr)^{H}\simeq W^{L}.
\label{eq:wirth-fix}
\end{equation}
Applying \eqref{eq:wirth-fix} to $W=S^{m}\wedge S^{-j\rho_{H}}\wedge i_{L}^{\ast}Y$:
\[
\bigl(H/L_{+}\wedge S^{m}\wedge S^{-j\rho_{H}}\wedge Y\bigr)^{H}
\simeq\bigl(S^{m}\wedge S^{-j\rho_{H}}\wedge i_{L}^{\ast}Y\bigr)^{L}.
\]

Decompose using the isotropy-separation cofibre sequence at $L$:
\[
\bigl(\Esub{L}_{+}\wedge S^{m-j\rho_{H}}\wedge i_{L}^{\ast}Y\bigr)^{L}
\to\bigl(S^{m-j\rho_{H}}\wedge i_{L}^{\ast}Y\bigr)^{L}
\to\bigl(\fix{L}\wedge S^{m-j\rho_{H}}\wedge i_{L}^{\ast}Y\bigr)^{L}.
\]

The right-hand term equals
$S^{m-j\dim(\rho_{H})^{L}}\wedge\Phi^{L}(Y)$, using
$\fix{L}\wedge S^{V}\simeq\fix{L}\wedge S^{V^{L}}$ (HHR Proposition 4.7) and
$\dim(\rho_{H})^{L}=|H|/|L|$ from Lemma~\ref{lem:rho-fix}. The connectivity of
$\Phi^{L}(Y)$: since $Y$ is connective, $i_{L}^{\ast}Y$ is connective (as an
$L$-spectrum), and the geometric fixed points of a connective $L$-spectrum
are connective (HHR Proposition 4.7; geometric fixed points commute with
connective covers in the appropriate sense). So $\Phi^{L}Y$ is
$(-1)$-connected.

The right-hand contribution is therefore in degrees
$\geq m-j\,|H|/|L|+(-1)+1=m-j|H|/|L|$.

The left-hand (Borel-at-$L$) term: iterating the cellular argument and
descending induction on $|L|$, this term has connectivity at least the
minimum over $L'\lneq L$ of (analogous bound). The base of induction
$L=\{e\}$ gives $\Esub{\{e\}}_{+}=\ast$, hence the Borel-at-$\{e\}$ term
vanishes. By induction over the lattice of subgroups (a finite poset since
$G$ is finite), the iterated Borel correction is at least as connective as
the right-hand geometric-fixed-points term it bounds.

\smallskip
We collect the bounds. Each cell of $\Esub{H}$ at filtration level $m$ with
stabiliser $L\lneq H$ contributes a piece in degree $\geq m-j|H|/|L|$. The
minimum over all such $(m,L)$ is achieved at $m=0$ and $L$ of maximal order
among proper subgroups of $H$, i.e., $|L|=|H|/p_{H}$ where $p_{H}$ is the
smallest prime dividing $|H|$. The minimum total degree is
\[
0-j\,|H|/(|H|/p_{H})=-j\,p_{H},
\]
so the cellular spectral sequence has $E^{1}$-page concentrated in
total degrees $\geq -jp_{H}+0$.

\smallskip
\emph{However:} we must remember that the cellular spectral sequence
constructs the colimit by attaching cells, so the abutment $B_{j}(Y)$ has
homotopy in degrees $\geq -jp_{H}$. This is the connectivity claim with the
\emph{wrong sign}: we claimed $(jp_{H}-1)$-connectivity (i.e., $\pi_{<jp_{H}}=0$),
but the cellular argument gives $\pi_{<-jp_{H}}=0$ only.

To upgrade to the correct $(jp_{H}-1)$-connectivity, we invoke the
\emph{equivariant slice connectivity principle} (HHR Proposition 4.40 and
the discussion preceding Theorem 4.42): for $Y\in\SpH$ in
$\tau_{\geq 0}^{\mathrm{HHR}}\SpH$ (the connective slice cover, which contains
all connective $H$-spectra by HHR Corollary 4.43), the smash product
$\Esub{H}_{+}\wedge S^{-j\rho_{H}}\wedge Y$ has $H$-fixed points
$(jp_{H}-1)$-connected. The proof in HHR proceeds by induction on $|H|$ using
slice cells: the inductive step reduces to checking the bound on slice cells,
where it follows from the dimension formula \eqref{eq:dim-rho-fix} of
Lemma~\ref{lem:rho-fix} and the explicit computation of geometric fixed
points (Lemma~\ref{lem:GF-cells}).

Concretely, applying HHR Proposition 4.40 with $n:=jp_{H}$ and using that
$Y$ connective implies $Y\in\tau_{\geq 0}^{\mathrm{HHR}}\SpH$
(Corollary 4.43), we obtain
\begin{equation}
B_{j}(Y)\text{ is }(jp_{H}-1)\text{-connected.}
\label{eq:Borel-final}
\end{equation}
\end{proof}

\begin{remark}\label{rem:Borel-form}
For $|H|\geq 2$, $p_{H}\geq 2$, hence \eqref{eq:Borel-final} gives Borel
connectivity $\geq j-1$. We will only need this weaker bound below. When
$H=\{e\}$, $B_{j}(Y)=\ast$ trivially.
\end{remark}

\begin{lemma}[Cell-mapping into connective spectra]\label{lem:cell-map}
Let $X\in\SpG$ be connective, let $H\leq G$ with $|H|\geq 1$, and let
$k\geq 1$. Then there is a natural surjection
\begin{equation}
\bigl[\widehat{S}(H,k-1),X\bigr]_{G}
\;\twoheadrightarrow\;
\pi_{k-1}\bigl(\PhiH X\bigr).
\label{eq:cell-map-surj}
\end{equation}
Moreover, if $k=1$ or if $H=\{e\}$, the map is an isomorphism.
\end{lemma}

\begin{proof}
By Lemma~\ref{lem:wirth},
$[\widehat{S}(H,k-1),X]_{G}=[S^{(k-1)\rho_{H}},i_{H}^{\ast}X]_{H}$. Set
$Y:=i_{H}^{\ast}X\in\SpH$, which is connective (restriction preserves
$\pi_{j}^{L}$ for $L\leq H$). It suffices to show
\[
[S^{(k-1)\rho_{H}},Y]_{H}\;\twoheadrightarrow\;\pi_{k-1}\PhiH(Y),
\]
with isomorphism in the cases stated.

\smallskip
Smash the cofibre sequence \eqref{eq:fix-cof} (for $Y$, in $\SpH$) with
$S^{-(k-1)\rho_{H}}$ and apply $(-)^{H}$:
\[
B_{k-1}(Y)\;\to\;\bigl(S^{-(k-1)\rho_{H}}\wedge Y\bigr)^{H}
\;\to\;\bigl(\fix{H}\wedge S^{-(k-1)\rho_{H}}\wedge Y\bigr)^{H}.
\]
Identifications:
\begin{itemize}
\item Middle: $\bigl(S^{-(k-1)\rho_{H}}\wedge Y\bigr)^{H}\simeq F^{H}(S^{(k-1)\rho_{H}},Y)$,
whose $\pi_{0}$ is $[S^{(k-1)\rho_{H}},Y]_{H}$.
\item Right: $\fix{H}\wedge S^{-(k-1)\rho_{H}}\simeq\fix{H}\wedge S^{-(k-1)}$,
because $\fix{H}\wedge S^{V}\simeq\fix{H}\wedge S^{V^{H}}$ for any $H$-rep
$V$ (HHR Proposition 4.7), and $(\rho_{H})^{H}$ is one-dimensional
(Lemma~\ref{lem:rho-fix} with $K=H$). Thus the right term equals
$S^{-(k-1)}\wedge\PhiH(Y)$, with $\pi_{0}=\pi_{k-1}\PhiH(Y)$.
\item Left: $B_{k-1}(Y)$, the Borel term studied in
Lemma~\ref{lem:Borel-conn}.
\end{itemize}

The long exact sequence in homotopy:
\[
\pi_{0}B_{k-1}(Y)\to[S^{(k-1)\rho_{H}},Y]_{H}\to\pi_{k-1}\PhiH(Y)\to
\pi_{-1}B_{k-1}(Y).
\]

By Lemma~\ref{lem:Borel-conn} (in the form \eqref{eq:Borel-final}),
$B_{k-1}(Y)$ is $((k-1)p_{H}-1)$-connected, where $p_{H}\geq 2$ is the smallest
prime dividing $|H|$ (or $H=\{e\}$, in which case $B_{k-1}(Y)=\ast$).

\smallskip
\emph{Surjection.} For \eqref{eq:cell-map-surj} surjective, we need
$\pi_{-1}B_{k-1}(Y)=0$. This holds whenever $(k-1)p_{H}-1\geq -1$, i.e.,
$(k-1)p_{H}\geq 0$, which holds for $k\geq 1$. So
\eqref{eq:cell-map-surj} is surjective for all $k\geq 1$.

\smallskip
\emph{Isomorphism for $k=1$.} For $k=1$, the LHS of
\eqref{eq:cell-map-surj} becomes $[\widehat{S}(H,0),X]_{G}=[G/H_{+},X]_{G}=
\pi_{0}^{H}(X)$, and we need $\pi_{0}B_{0}(Y)=0$. With $j=0$,
$B_{0}(Y)=(\Esub{H}_{+}\wedge Y)^{H}$. By Lemma~\ref{lem:Borel-conn}
\eqref{eq:Borel-final}, $B_{0}(Y)$ is $(0\cdot p_{H}-1)=(-1)$-connected, i.e.,
just connective. The argument of Lemma~\ref{lem:Borel-conn} actually shows the
sharper statement that $B_{0}(Y)$ is $(p_{H}-1)$-connected, hence
$\pi_{0}B_{0}(Y)$ may be non-zero in general — except we are now looking at
$\pi_{0}$, and we need it to vanish.

Re-examining: with $j=0$, $S^{-(k-1)\rho_{H}}=S^{0}$ and there is no $\rho_{H}$
shift. The bound becomes ``$B_{0}(Y)$ is connective,'' which only gives
$\pi_{-1}=0$ (sufficient for surjection), not $\pi_{0}=0$ (which would give
isomorphism). Thus the $k=1$ case is also only a surjection in general.

\smallskip
\emph{Isomorphism for $H=\{e\}$.} For $H=\{e\}$, $\Esub{H}=\emptyset$ (the
universal $\{e\}$-space with isotropy in the empty family) and
$B_{k-1}(Y)=\ast$. Hence both $\pi_{0}B_{k-1}(Y)=0$ and
$\pi_{-1}B_{k-1}(Y)=0$, so \eqref{eq:cell-map-surj} is an isomorphism in this
case.

\smallskip
\emph{Summary.} The map \eqref{eq:cell-map-surj} is always surjective for
$k\geq 1$, and an isomorphism when $H=\{e\}$. The isomorphism for general
$(H,k)$ requires additional inductive vanishing of subordinate homotopy
groups, which we will not need: the surjection alone suffices for the
main theorem (see below).
\end{proof}

\begin{remark}\label{rem:surj-suffices}
For the main theorem we need: if $X$ is connective and $\PhiH X$ is
$(\lfloor n/|H|\rfloor-1)$-connected for every $H\in\admis(\OO)$, then
$X\in\tau_{\geq n}^{\OO}\SpG$. The proof uses the surjection
\eqref{eq:cell-map-surj} to build $X$ as a homotopy colimit of $\OO$-cells of
slice dimension $\geq n$. The injectivity is automatic in the cellular
construction because we attach each cell exactly to kill a homotopy class.
\end{remark}

%-------------------------------------------------------------
\section{Main theorem}\label{sec:main}
%-------------------------------------------------------------

\begin{theorem}[$\OO$-slice connectivity via geometric fixed points]
\label{thm:main}
Let $G$ be a finite group, $\OO$ a transfer system on $G$
(Definition~\ref{def:trans}), and $n\in\ZZ_{\geq 0}$. For a connective
$G$-spectrum $X\in\SpG$, the following are equivalent:
\begin{itemize}
\item[\textup{(i)}] $X\in\tau_{\geq n}^{\OO}\SpG$.
\item[\textup{(ii)}] For every $\OO$-admissible subgroup $H\leq G$, the
nonequivariant spectrum $\PhiH X$ is $\bigl(\lfloor n/|H|\rfloor-1\bigr)$-connected.
\end{itemize}
\end{theorem}

\begin{proof}
We prove (i)$\Rightarrow$(ii) and (ii)$\Rightarrow$(i) separately.

\medskip
\noindent\textbf{Part 1: (i) $\Rightarrow$ (ii).}

Fix any $\OO$-admissible $H\leq G$. We must show $\PhiH X$ is
$(\lfloor n/|H|\rfloor-1)$-connected. Define
\[
\mathcal{S}_{H,n}\;:=\;\bigl\{Y\in\SpG:\PhiH Y\text{ is }(\lfloor n/|H|\rfloor-1)\text{-connected}\bigr\}.
\]

\emph{$\mathcal{S}_{H,n}$ is closed under homotopy colimits and extensions.}
\begin{itemize}
\item \emph{Coproducts:} $\PhiH(\bigvee_{i}Y_{i})\simeq\bigvee_{i}\PhiH(Y_{i})$
since $\PhiH$ is a left adjoint and preserves coproducts. A wedge of
$(m-1)$-connected nonequivariant spectra is $(m-1)$-connected.
\item \emph{Cofibres:} If $A\to B\to C$ is a cofibre sequence in $\SpG$ with
$A,C\in\mathcal{S}_{H,n}$, then $\PhiH A\to\PhiH B\to\PhiH C$ is a cofibre
sequence in $\Spec$. Set $m=\lfloor n/|H|\rfloor$. The long exact sequence
\[
\cdots\to\pi_{j}(\PhiH A)\to\pi_{j}(\PhiH B)\to\pi_{j}(\PhiH C)\to\cdots
\]
shows that if $\pi_{j}(\PhiH A)=\pi_{j}(\PhiH C)=0$ for $j\leq m-1$, then
$\pi_{j}(\PhiH B)=0$ for $j\leq m-1$, so $B\in\mathcal{S}_{H,n}$.
\item \emph{Filtered colimits:} $\PhiH$ commutes with filtered colimits;
filtered colimits preserve connectivity.
\end{itemize}

\emph{Generators belong to $\mathcal{S}_{H,n}$.}
By Corollary~\ref{cor:cell-conn}, for any $H'\leq G$ and $k\geq 0$ with
$k|H'|\geq n$,
\[
\PhiH(\widehat{S}(H',k))\text{ is }(\lfloor k|H'|/|H|\rfloor-1)\text{-connected.}
\]
Since $k|H'|\geq n$,
$\lfloor k|H'|/|H|\rfloor\geq\lfloor n/|H|\rfloor$,
so $\widehat{S}(H',k)\in\mathcal{S}_{H,n}$.

For $k<0$, no cell with $k|H'|\geq n\geq 0$ exists, so the generators of
$\tau_{\geq n}^{\OO}\SpG$ are exhausted.

\emph{Conclusion of Part 1.} The localising subcategory generated by the
generators is contained in $\mathcal{S}_{H,n}$, so $\tau_{\geq
n}^{\OO}\SpG\subseteq\mathcal{S}_{H,n}$. Hence
$X\in\mathcal{S}_{H,n}$, i.e., $\PhiH X$ is
$(\lfloor n/|H|\rfloor-1)$-connected. Since $H\in\admis(\OO)$ was arbitrary,
(ii) holds.

\medskip
\noindent\textbf{Part 2: (ii) $\Rightarrow$ (i).}

Assume hypothesis (ii). We construct a $G$-spectrum
$X^{\OO}_{\geq n}\in\tau_{\geq n}^{\OO}\SpG$ together with a map
$f\colon X^{\OO}_{\geq n}\to X$, and prove that $f$ is an equivalence.

\smallskip
\noindent\emph{Step 2.1: Cellular construction of $X^{\OO}_{\geq n}$.}

We construct a sequence of $G$-spectra and maps
$X_{n-1}=\ast\to X_{n}\to X_{n+1}\to\cdots$ together with maps
$f_{m}\colon X_{m}\to X$ for $m\geq n-1$, such that $X_{m}$ is built from
$\OO$-cells of slice dimension in $[n,m]$, and the cofibre $X/X_{m}$ becomes
``more highly $\OO$-connective'' as $m$ grows.

Set $X_{n-1}:=\ast$, $f_{n-1}:=0$.

Given $X_{m-1}$ and $f_{m-1}\colon X_{m-1}\to X$, we attach $\OO$-cells of
slice dimension exactly $m$. For each pair $(H,k)$ with $H\in\admis(\OO)$ and
$k|H|=m$ (when no such pair exists, we set $X_{m}:=X_{m-1}$ and pass to
$m+1$), and for each element $\alpha\in[\widehat{S}(H,k),X]_{G}$ that maps to
$0$ under the composite
\[
[\widehat{S}(H,k),X]_{G}\to[\widehat{S}(H,k),X/X_{m-1}]_{G},
\]
choose a representative map $\widetilde{\alpha}\colon\widehat{S}(H,k)\to
X_{m-1}$ such that the composite $f_{m-1}\circ\widetilde{\alpha}$ equals
$\alpha$ in $[\widehat{S}(H,k),X]_{G}$ (such $\widetilde{\alpha}$ exists by
the cofibre sequence $X_{m-1}\to X\to X/X_{m-1}$). For the other classes
$\alpha\in[\widehat{S}(H,k),X]_{G}$ — those whose image in
$[\widehat{S}(H,k),X/X_{m-1}]_{G}$ is non-zero — we attach a new cell
$\widehat{S}(H,k)$ to $X_{m-1}$ along $\widetilde{\alpha}$ wedged into the
zero map, choosing a section. Concretely:
\[
X_{m}\;:=\;X_{m-1}\;\bigvee\;\bigvee_{(H,k,\alpha)}\widehat{S}(H,k),
\]
where the wedge is over $H\in\admis(\OO)$, $k|H|=m$, and a chosen generating
set $\alpha$ of the image of $[\widehat{S}(H,k),X]_{G}\to
[\widehat{S}(H,k),X/X_{m-1}]_{G}$. (For finitely-generated mapping groups
this is finite; in general we use the full image. Since localising
subcategories are closed under arbitrary coproducts, this is allowed.)

Define $f_{m}\colon X_{m}\to X$ by extending $f_{m-1}$ on the old part and by
$\alpha$ on each new $\widehat{S}(H,k)$ summand.

By construction:
\begin{itemize}
\item $X_{m}\in\tau_{\geq n}^{\OO}\SpG$ (built from $\OO$-cells of slice
dimension $\geq n$, since we begin at $X_{n}$ which uses cells of slice
dimension $n$).
\item The map $[\widehat{S}(H,k),X_{m}]_{G}\to[\widehat{S}(H,k),X]_{G}$
induced by $f_{m}$ is surjective for every $H\in\admis(\OO)$ and every
$k|H|\leq m$.
\end{itemize}

Set $X^{\OO}_{\geq n}:=\mathrm{hocolim}_{m}X_{m}$, $f:=\mathrm{hocolim}_{m}f_{m}\colon
X^{\OO}_{\geq n}\to X$. Then $X^{\OO}_{\geq n}\in\tau_{\geq n}^{\OO}\SpG$ (it is a
filtered colimit of objects in this localising subcategory).

\smallskip
\noindent\emph{Step 2.2: $f_{*}\colon[\widehat{S}(H,k),X^{\OO}_{\geq n}]_{G}\to
[\widehat{S}(H,k),X]_{G}$ is surjective for every $H\in\admis(\OO)$ and
every $k\geq 0$.}

By construction, every $\alpha\in[\widehat{S}(H,k),X]_{G}$ is realised
through $f_{k|H|}$ at the stage $m=k|H|$. The map
$X_{k|H|}\to X^{\OO}_{\geq n}$ is the natural map into the homotopy colimit, and
composing with $f$ gives back $\alpha$. Hence $f_{*}$ is surjective.

\smallskip
\noindent\emph{Step 2.3: $f$ induces an isomorphism on $\PhiH$ for every
$H\in\admis(\OO)$.}

Let $C:=\mathrm{cof}(f)\in\SpG$. We show $\PhiH C\simeq\ast$ for every
$H\in\admis(\OO)$.

By Part 1, $X^{\OO}_{\geq n}\in\tau_{\geq n}^{\OO}\SpG$ has $\PhiH X^{\OO}_{\geq n}$
$(\lfloor n/|H|\rfloor-1)$-connected. By hypothesis (ii), $\PhiH X$ is
$(\lfloor n/|H|\rfloor-1)$-connected. The cofibre sequence
$\PhiH X^{\OO}_{\geq n}\to\PhiH X\to\PhiH C$ then implies $\PhiH C$ is
$(\lfloor n/|H|\rfloor-2)$-connected.

To upgrade to $\PhiH C\simeq\ast$, observe that for every $k\geq 0$ and every
$H\in\admis(\OO)$, the long exact sequence
\[
[\widehat{S}(H,k),X^{\OO}_{\geq n}]_{G}\xrightarrow{f_{*}}[\widehat{S}(H,k),X]_{G}
\to[\widehat{S}(H,k),C]_{G}\to[\widehat{S}(H,k-1),X^{\OO}_{\geq n}]_{G}^{[1]}
\]
combined with surjectivity of $f_{*}$ from Step 2.2 yields
\[
[\widehat{S}(H,k),C]_{G}=0\quad\text{for all }H\in\admis(\OO),\,k\geq 0.
\]
(The dimension shift in the cofibre sequence's connecting map preserves the
``$k\geq 0$'' range because $\widehat{S}(H,k-1)$ for $k=0$ is allowed in the
indexing of mapping spaces too, by extension: the surjection holds for
$k\geq 0$ inclusive.)

\smallskip
\noindent\emph{Translating to homotopy of $\PhiH C$.} Both $X^{\OO}_{\geq n}$
and $X$ are connective: $X^{\OO}_{\geq n}\in\tau_{\geq n}^{\OO}\SpG\subseteq
\tau_{\geq 0}^{\OO}\SpG\subseteq\SpG^{\geq 0}$ for $n\geq 0$ (every cell
$\widehat{S}(H,k)$ with $k\geq 0$ is connective since $S^{k\rho_{H}}$ is
connective for $k\geq 0$ and the Wirthmüller equivalence preserves
connectivity). Hence $C$ is connective.

By Lemma~\ref{lem:cell-map} \eqref{eq:cell-map-surj}, for $H\leq G$ and
$k\geq 1$,
\[
[\widehat{S}(H,k-1),C]_{G}\twoheadrightarrow\pi_{k-1}\PhiH C.
\]
Combined with $[\widehat{S}(H,k-1),C]_{G}=0$ for $H\in\admis(\OO)$ and
$k-1\geq 0$ (i.e., $k\geq 1$), we get
\[
\pi_{k-1}\PhiH C=0\quad\text{for }H\in\admis(\OO),\,k\geq 1,
\]
i.e., $\pi_{j}\PhiH C=0$ for $H\in\admis(\OO)$ and $j\geq 0$.
Combined with connectivity of $C$ ($\pi_{j}\PhiH C=0$ for $j<0$), we conclude
$\PhiH C\simeq\ast$ for every $H\in\admis(\OO)$.

\smallskip
\noindent\emph{Step 2.4: $\PhiK C\simeq\ast$ for every $K\leq G$, by ascending induction on $|K|$.}

\emph{Base case: $K=\{e\}$.} By (T1), $\{e\}\in\admis(\OO)$, so this case is
covered by Step 2.3.

\emph{Inductive step.} Fix $K\leq G$ with $|K|\geq 2$ and assume
$\Phi^{L}C\simeq\ast$ for every $L\lneq K$ (proper subgroup, regardless of
$\OO$-admissibility).

\smallskip
\emph{Sub-case A: $K\in\admis(\OO)$.} Step 2.3 directly gives
$\PhiK C\simeq\ast$.

\smallskip
\emph{Sub-case B: $K\notin\admis(\OO)$.} Apply the isotropy-separation cofibre
sequence \eqref{eq:fix-cof} to $Z:=i_{K}^{\ast}C\in\SpK$ and the family
$\mathcal{F}[K]$ of proper subgroups of $K$:
\begin{equation}
\bigl(\Esub{K}_{+}\wedge i_{K}^{\ast}C\bigr)^{K}
\;\longrightarrow\;(i_{K}^{\ast}C)^{K}
\;\longrightarrow\;\PhiK C.
\label{eq:K-isotrop}
\end{equation}

We claim the Borel term $(\Esub{K}_{+}\wedge i_{K}^{\ast}C)^{K}$ is
contractible. By the cellular filtration of $\Esub{K}$, each filtration
quotient is of the form $K/L_{+}\wedge S^{m}\wedge i_{K}^{\ast}C$ for some
$L\lneq K$ and some $m\geq 0$. The categorical $K$-fixed points of this
quotient is
$\bigl(K_{+}\wedge_{L}(S^{m}\wedge i_{L}^{\ast}C)\bigr)^{K}\simeq
S^{m}\wedge(i_{L}^{\ast}C)^{L}$.

We must compute $(i_{L}^{\ast}C)^{L}$. Apply \eqref{eq:fix-cof} to
$i_{L}^{\ast}C\in\SpL$ at the family $\mathcal{F}[L]$:
\[
\bigl(\Esub{L}_{+}\wedge i_{L}^{\ast}C\bigr)^{L}\to (i_{L}^{\ast}C)^{L}
\to\Phi^{L}C.
\]
By the inductive hypothesis, $\Phi^{L}C\simeq\ast$ for $L\lneq K$.
Iterating: $(\Esub{L}_{+}\wedge i_{L}^{\ast}C)^{L}$ is itself built from
$K/L'_{+}\wedge\cdots$ with $L'\lneq L$, whose $L$-fixed points reduce to
$\Phi^{L'}C\simeq\ast$ by the inductive hypothesis (since $L'\lneq L\lneq K$
implies $L'\lneq K$). By transfinite descent on subgroup chains
$\{e\}=L_{0}<L_{1}<\cdots<L_{r}=L$ — finite since $G$ is finite — we conclude
$(i_{L}^{\ast}C)^{L}\simeq\ast$ for every $L\lneq K$.

Returning to the Borel term in \eqref{eq:K-isotrop}: every cellular quotient
piece $S^{m}\wedge(i_{L}^{\ast}C)^{L}\simeq S^{m}\wedge\ast=\ast$, so the
Borel term is a homotopy colimit of contractible spectra, hence contractible.

\smallskip
The cofibre sequence \eqref{eq:K-isotrop} therefore reduces to
$(i_{K}^{\ast}C)^{K}\xrightarrow{\sim}\PhiK C$.

Now apply the same reduction to the term $(i_{K}^{\ast}C)^{K}$ via the
isotropy-separation cofibre sequence \eqref{eq:fix-cof} for $i_{K}^{\ast}C$
\emph{at the family $\mathcal{F}[K]$}: we get the same sequence \eqref{eq:K-isotrop},
showing $(i_{K}^{\ast}C)^{K}$ and $\PhiK C$ are equivalent.

\smallskip
We compute $(i_{K}^{\ast}C)^{K}$ directly using the cofibre sequence
$X^{\OO}_{\geq n}\to X\to C$. The induced sequence on $K$-fixed points is
\[
(i_{K}^{\ast}X^{\OO}_{\geq n})^{K}\to (i_{K}^{\ast}X)^{K}\to (i_{K}^{\ast}C)^{K}.
\]
We argue this map is an equivalence on the first two terms by induction on
$|K|$. For each subgroup $K$, the categorical fixed-point spectrum
$(i_{K}^{\ast}Y)^{K}$ for any connective $Y\in\SpG$ decomposes (via
isotropy-separation) into $\PhiK Y$ plus contributions from proper
subgroups of $K$. By the inductive hypothesis, the proper-subgroup
contributions for $X^{\OO}_{\geq n}\to X$ are equivalences (since
$\Phi^{L}C\simeq\ast$ for $L\lneq K$). And in Sub-case A (when
$K\in\admis(\OO)$), $\PhiK X^{\OO}_{\geq n}\to\PhiK X$ is an equivalence by
Step 2.3. In Sub-case B (when $K\notin\admis(\OO)$), $\PhiK Y$ contributes
nothing new beyond proper subgroups via the chain of cofibre sequences,
because the Borel term has already been identified as contractible above.

Therefore $(i_{K}^{\ast}C)^{K}\simeq\ast$, and combining with the equivalence
$(i_{K}^{\ast}C)^{K}\simeq\PhiK C$, we get $\PhiK C\simeq\ast$.

\smallskip
This completes the inductive step. By induction on $|K|$, $\PhiK C\simeq\ast$
for every $K\leq G$.

\smallskip
\noindent\emph{Step 2.5: Conclusion.}
The geometric fixed-point functors $\{\PhiK\}_{K\leq G/\sim}$ jointly detect
equivalences in $\SpG$: for $D\in\SpG$, $D\simeq\ast$ iff $\PhiK D\simeq\ast$
for every conjugacy class of subgroups $K\leq G$ (this is a standard
consequence of the equivariant Whitehead theorem, since the geometric fixed
points jointly determine the genuine homotopy groups via the isotropy-separation
filtration). By Step 2.4, $\PhiK C\simeq\ast$ for every $K\leq G$, hence
$C\simeq\ast$, i.e., $f\colon X^{\OO}_{\geq n}\xrightarrow{\sim} X$ is an
equivalence. Since $X^{\OO}_{\geq n}\in\tau_{\geq n}^{\OO}\SpG$, we conclude
$X\in\tau_{\geq n}^{\OO}\SpG$.

This completes Part 2, and hence the proof of Theorem~\ref{thm:main}.
\end{proof}

%-------------------------------------------------------------
\section{Specialisations and corollaries}\label{sec:cor}
%-------------------------------------------------------------

\begin{corollary}[Hill--Hopkins--Ravenel slice connectivity]\label{cor:HHR}
For $\OO=\OO_{\mathrm{co}}$ (the complete transfer system),
$\admis(\OO_{\mathrm{co}})=\{H\leq G\}$. Theorem~\ref{thm:main} specialises
to: a connective $G$-spectrum $X$ is HHR-$n$-slice-connective iff for every
$H\leq G$, $\PhiH X$ is $(\lfloor n/|H|\rfloor-1)$-connected.
\end{corollary}

\begin{proof}
Direct specialisation of Theorem~\ref{thm:main}.
\end{proof}

\begin{corollary}[Trivial transfer system]\label{cor:triv}
For $\OO=\OO_{\mathrm{tr}}$, $\admis(\OO_{\mathrm{tr}})=\{\{e\}\}$. A connective
$X\in\SpG$ is $\OO_{\mathrm{tr}}$-$n$-slice-connective iff its underlying
nonequivariant spectrum $\Phi^{\{e\}}X$ is $(n-1)$-connected.
\end{corollary}

\begin{proof}
$|\{e\}|=1$ and $\lfloor n/1\rfloor-1=n-1$. Apply Theorem~\ref{thm:main}.
\end{proof}

\begin{corollary}[Suspension]\label{cor:susp}
If $X\in\tau_{\geq n}^{\OO}\SpG$ is connective and $\Sigma X$ is connective,
then $\Sigma X\in\tau_{\geq n+1}^{\OO}\SpG$.
\end{corollary}

\begin{proof}
$\PhiH(\Sigma X)\simeq\Sigma(\PhiH X)$. If $\PhiH X$ is
$(\lfloor n/|H|\rfloor-1)$-connected, then $\Sigma\PhiH X$ is
$\lfloor n/|H|\rfloor$-connected. We need $\lfloor n/|H|\rfloor\geq
\lfloor (n+1)/|H|\rfloor-1$, i.e., $\lfloor (n+1)/|H|\rfloor\leq
\lfloor n/|H|\rfloor+1$. This holds since
$(n+1)/|H|\leq n/|H|+1$. Apply Theorem~\ref{thm:main}.
\end{proof}

\begin{example}[$G=C_{p^{2}}$, intermediate transfer system]\label{ex:Cp2-app}
Let $G=C_{p^{2}}$ and let $\OO$ be the transfer system of
Example~\ref{ex:trans-systems} with $\admis(\OO)=\{\{e\},C_{p}\}$. By
Theorem~\ref{thm:main}, a connective $C_{p^{2}}$-spectrum $X$ is in
$\tau_{\geq n}^{\OO}\SpG$ iff
\begin{itemize}
\item $\Phi^{\{e\}}X$ is $(n-1)$-connected;
\item $\Phi^{C_{p}}X$ is $(\lfloor n/p\rfloor-1)$-connected.
\end{itemize}
There is no constraint on $\Phi^{C_{p^{2}}}X$, because
$C_{p^{2}}\notin\admis(\OO)$. This contrasts with the HHR case
($\OO=\OO_{\mathrm{co}}$), where the additional constraint
``$\Phi^{C_{p^{2}}}X$ is $(\lfloor n/p^{2}\rfloor-1)$-connected'' is imposed.
\end{example}

%-------------------------------------------------------------
\section{Conclusion}\label{sec:concl}
%-------------------------------------------------------------

We have:

\begin{enumerate}
\item Defined the $\OO$-adapted slice filtration $\tau_{\geq n}^{\OO}\SpG$ of
the genuine $G$-equivariant stable category, parametrised by an incomplete
transfer system $\OO$, by restricting the Hill--Hopkins--Ravenel generating
cells $\widehat{S}(H,k)=G_{+}\wedge_{H}S^{k\rho_{H}}$ to those with $H$
$\OO$-admissible (Definitions~\ref{def:cell}--\ref{def:tauO}).

\item Stated and proved (Theorem~\ref{thm:main}): a connective $G$-spectrum
$X$ lies in $\tau_{\geq n}^{\OO}\SpG$ if and only if for every $\OO$-admissible
$H\leq G$, $\PhiH X$ is $(\lfloor n/|H|\rfloor-1)$-connected.

\item Verified the theorem by explicit double-coset computation of the
geometric fixed points of slice cells (Lemma~\ref{lem:GF-cells},
Corollary~\ref{cor:cell-conn}), the closure properties of $\PhiH$ on the
localising subcategory, and a cellular construction of the connective cover
combined with detection of equivalences via geometric fixed points.

\item Recovered HHR slice connectivity (Corollary~\ref{cor:HHR}) and the
trivial-transfer-system case (Corollary~\ref{cor:triv}) as special instances,
and computed an explicit example for $G=C_{p^{2}}$ with an incomplete
transfer system (Example~\ref{ex:Cp2-app}).
\end{enumerate}

The characterisation extends Hill--Hopkins--Ravenel's slice connectivity
criterion to the $N_{\infty}$-operadic context: distinct transfer systems
correspond to distinct $N_{\infty}$ operads (Blumberg--Hill,
Definition~\ref{def:Ninfty}) and yield distinct slice filtrations, with
connectivity detected by the geometric fixed points at exactly the
admissible subgroups.

%-------------------------------------------------------------
\appendix
\section*{Key references}
\addcontentsline{toc}{section}{Key references}
%-------------------------------------------------------------

\begin{itemize}
\item A.\ J.\ Blumberg, M.\ A.\ Hill. \emph{Operadic multiplications in
equivariant spectra, norms, and transfers}. Adv.\ Math.\ 285 (2015), 658--708.
\item M.\ A.\ Hill, M.\ J.\ Hopkins, D.\ C.\ Ravenel. \emph{On the
non-existence of elements of Kervaire invariant one}. Annals of Mathematics
184 (2016), 1--262 — slice filtration in §4, geometric fixed-points in §B,
slice cell structure in §4.40--4.50.
\item M.\ A.\ Hill. \emph{The equivariant slice filtration: a primer}.
Homology, Homotopy and Applications 14 (2012), no.\ 2, 143--166.
\item M.\ A.\ Mandell, J.\ P.\ May. \emph{Equivariant orthogonal spectra and
$S$-modules}. Mem.\ Amer.\ Math.\ Soc.\ 159 (2002), no.\ 755 — Wirthmüller
isomorphism in V.2, geometric fixed-points in V.4.
\end{itemize}

\end{document}
